\chapter{Machine Learning Methodology}
\section{Feature Matrix and Labels}
Let $X \in \mathbb{R}^{n \times p}$ be the final expression matrix, where $n$ is the number of single cells and $p=3000$ is the number of selected highly variable genes. The label vector $\mathbf{y}$ contains disease conditions from \texttt{condition.l1}.

\section{Random Forest Classifier}
Random Forest is an ensemble of decision trees. Each tree is trained on a bootstrap sample, and each split considers a random subset of features. For a cell $x$, the predicted class is obtained by majority voting:
\begin{equation}
\hat{y}=\operatorname{mode}\{h_1(x),h_2(x),\ldots,h_T(x)\},
\end{equation}
where $T=300$ is the number of trees and $h_t$ is the prediction of tree $t$.

The model was configured as:
\begin{equation}
\texttt{RandomForestClassifier}(n\_estimators=300, class\_weight=\text{balanced}, random\_state=42).
\end{equation}

\begin{figure}[H]
\centering
\includegraphics[width=0.9\textwidth]{random_forest_workflow.pdf}
\caption{Random Forest workflow for training multiple decision trees and aggregating predictions through majority voting.}
\label{fig:rf_workflow}
\end{figure}

\section{Feature Importance}
For gene $j$, Random Forest importance can be estimated by the average impurity reduction over all trees:
\begin{equation}
I_j = \frac{1}{T}\sum_{t=1}^{T}\sum_{s \in S_{j,t}} \Delta G(s),
\end{equation}
where $S_{j,t}$ is the set of splits in tree $t$ using gene $j$ and $\Delta G(s)$ is the decrease in Gini impurity at split $s$.

\begin{figure}[H]
\centering
\includegraphics[width=0.86\textwidth]{feature_importance_plot.pdf}
\caption{Top genes ranked by Random Forest feature importance. Highly ranked genes contribute strongly to disease-state classification.}
\label{fig:feature_importance}
\end{figure}

\section{Evaluation Metrics}
For each class, precision, recall, and F1-score are defined as
\begin{equation}
Precision=\frac{TP}{TP+FP}, \qquad Recall=\frac{TP}{TP+FN},
\end{equation}
\begin{equation}
F1 = 2\times\frac{Precision\times Recall}{Precision+Recall}.
\end{equation}
The confusion matrix $C$ is defined by
\begin{equation}
C_{ij}=|\{k: y_k=i \land \hat{y}_k=j\}|,
\end{equation}
where rows represent true classes and columns represent predicted classes.
\begin{table}[H]
\centering
\caption{Random Forest evaluation metrics template. Values should be filled from \texttt{results/RF\_classification\_report.csv}.}
\label{tab:rf_metrics}
\begin{tabular}{lrrrr}
\toprule
Class & Precision & Recall & F1-score & Support \\
\midrule
Ref & -- & -- & -- & -- \\
AKI & -- & -- & -- & -- \\
CKD & -- & -- & -- & -- \\
Macro average & -- & -- & -- & -- \\
Weighted average & -- & -- & -- & -- \\
\bottomrule
\end{tabular}
\end{table}


\begin{figure}[H]
\centering
\includegraphics[width=0.74\textwidth]{confusion_matrix.pdf}
\caption{Confusion matrix summarizing agreement between true and predicted disease-state labels.}
\label{fig:confusion_matrix}
\end{figure}

\begin{figure}[H]
\centering
\includegraphics[width=0.74\textwidth]{roc_curve.pdf}
\caption{ROC curve for pairwise or one-versus-rest disease-state classification. Higher area under the curve indicates stronger discriminative performance.}
\label{fig:roc_curve}
\end{figure}

\section{Biomarker Ranking}
For each pairwise comparison, candidate biomarkers were ranked using
\begin{equation}
Score_j = I_j \times |\bar{x}_{j,A}-\bar{x}_{j,B}|,
\end{equation}
where $I_j$ is the Random Forest importance of gene $j$, and $|\bar{x}_{j,A}-\bar{x}_{j,B}|$ is the absolute expression difference between conditions $A$ and $B$. This score favors genes with both predictive importance and measurable expression shifts.
